\section{A separate Grothendieck dictionary: objects, operations, and proof transport}
\label{sec:grothendieck-dictionary}

This part restates the checked deductions in the requested geometric language. A locator identifies the framework in which a proof lives; it is not a claim that Grothendieck stated a modern complexity result. The original manuscript images are not in the attachment. Cote and scan references below therefore refer to the supplied transcriptions, and the mathematical corrections do not certify an unseen handwritten reading.

\subsection{Exact framework locators and scope}

\paragraph{D-S3: group schemes and weights.}
For an affine group scheme, its representations are comodules over its Hopf algebra. The relevant reference is SGA~3, Expos\'e~I, Proposition~4.7.2. For a diagonalizable group $D(M)$, such comodules are $M$-graded modules, by Proposition~4.7.3; the following remark gives the corresponding graded-algebra description of affine schemes with action. Proposition~4.7.4 gives equivariant splittings when the underlying sequence splits. These statements, inspected in the original Expos\'e~I PDF, give the integral meaning of every torus weight and cocharacter below. For stabilizers, SGA~3, Expos\'e~VI$_B$, Proposition~4.1 gives upper semicontinuity of group-fiber dimension and explicitly cites EGA~IV$_3$, Theorem~13.1.3. The latter locator is corroborated through that primary SGA text; its entire EGA proof has not been reread here.

\paragraph{D-S6: $K_0$, exterior powers, Adams operations, and filtration.}
SGA~6, Expos\'e~V, Definition~2.1 and Example~2.2(b), construct the pre-$\lambda$ structure of locally projective modules over a ringed topos; Lemma~2.2.1 uses a filtration on exterior powers of an exact sequence. Definition~2.4 and equation~(2.4.2) specify the universal product and composition identities. Expos\'e~VI, section~3.2 gives the special-$\lambda$-ring assertion for the $K'$ of a quasi-compact scheme. Expos\'e~V, Proposition~3.2 defines $\gamma$ from $\lambda$; section~3.10, equation~(3.10.1), defines the augmentation filtration. Section~7.1, equation~(7.1.1), and Proposition~7.5 give Adams operations and their ring identities. The electronic transcriptions of the cited Expos\'e~V passages were inspected; the exact VI~3.2 statement was corroborated in the primary-text transcription during independent review. We do not describe an integral character identity as an integral semisimple decomposition.

For split reductive group schemes over a principal ideal domain, in particular over $\mathbb Z$, a useful complementary reference is Serre (1968), Theorem~5, section~3.7, and the explicit $\mathrm{GL}_n$ character ring in section~3.8. The base-change statements in that theorem concern the representation Grothendieck group, not preservation of a characteristic-zero basis of simple objects in every characteristic. Nor do they automatically answer the separate finite-symmetric-group lattice question recorded in the supplied study of cote~133.


\paragraph{D-Tor: derived restriction and its obstruction groups.}
SGA~6, Expos\'e~III, Definition~1.5 defines Tor independence of two
schemes over a base by vanishing of their positive Tor sheaves; in
the affine case it is exactly the assertion that derived tensor
product agrees with ordinary tensor product. Notation~1.6 specifies
the relative derived tensor product. Expos\'e~I, Proposition~5.1
and Definition~5.2 give the flat-amplitude criterion, and paragraph~5.4
identifies it with ordinary finite flat-resolution length for a
module. These exact passages were read in the electronic French
primary-text transcription. III~1.5 refers to EGA~III, section~6.5,
for local hyper-Tor; that reference is corroborated here through
SGA~6, without a new whole-section reading of EGA~III.

\paragraph{D-FGA: images, quotients, and projective parameter spaces.}
For a homogeneous orbit map over $\mathbb Z$, take its schematic image by the kernel of the coordinate homomorphism, as in G-04. Degreewise kernels and images are given by finite matrices and their minors. EGA~II, sections~2.4--2.5, in particular Proposition~2.5.2, gives the projective scheme and associated sheaf of a graded module. The electronic EGA translation was inspected for this locator. FGA, Bourbaki Expos\'e~221, Theorem~3.1 gives the Quot parameter space, Theorem~3.2 its fixed-Hilbert-polynomial projective form, and Proposition~3.8 its Grassmannian realization in the required degree range. The original FGA scan was inspected. These results place the families in parameter spaces; they supply no polynomial bound for finding the relevant minors or their nonvanishing loci.

\paragraph{D-ET: traces, weights, and their arithmetic size.}
SGA~4$\frac12$, \emph{Rapport sur la formule des traces}, Theorems~3.1--3.2 and Proposition~3.3, gives the determinant expression for the $L$-function, the trace formula, and the logarithmic-derivative identity. SGA~5, Expos\'e~X, Theorem~7.1 and Corollary~7.12, formula~(7.13), gives the Euler--Poincar\'e conductor accounting. Deligne's Weil~II, Theorem~3.3.1 and section~3.3.10, supplies the compact-support weight estimate used in P-06; Theorem~3.5.3 gives the equidistribution setting. The Weil~II locators were checked against Katz and Sarnak (1999), sections~9.0.13--9.0.16 and Theorem~9.2.6. The original large Weil~II PDF was not successfully opened here. This is a statement-and-dependency check, not an independent verification of the complete proof of Weil~II.

\subsection{Rectangular characters and orbit-closure multiplicities}

\paragraph{R-01. Tensor pullback of a Schur operation.}
For locally free $E,F$ of ranks $a,n$, the tensor morphism $\mathrm{GL}(E)\times\mathrm{GL}(F)\to\mathrm{GL}(E\otimes F)$ is defined over $\mathbb Z$. The class $s_{(d^n)}[E\otimes F]$ restricts, after passage to characteristic zero, to the sum of $[S_\mu E\otimes S_\nu F]$ with coefficients $g((d^n),\mu,\nu)$. The Haar projector is the complex realization of extraction of that character component. Its finite CUE dimension is a rank of a Schur construction; its normalized limiting constant need not be an integer rank. This is precisely the proof of R-01 transported through D-S3 and D-S6. Cote~133, scan~028, lines~6--17, provides the tensor-power and symmetric-operation language, while scan~019 distinguishes rank from the full character class.

\paragraph{R-02. Two different meanings of positivity.}
The tableau expansion is an integral character identity in the universal special $\lambda$-ring, with the Weyl alternating factor extracting a highest weight. Over $\mathbb C$, the invariant summand of $S_\lambda(E\otimes F)\otimes\det(E)^{-d}\otimes\det(F)^{-d}$ is an actual vector space. A Hermitian metric identifies its projector trace with a sum of squared absolute values. The former is an integral identity; the latter uses complex conjugation and compact integration. D-S6 and cote~133, scan~028, lines~19--24 (the direct-summand question), are the exact conceptual distinction. No cancellation-free combinatorial basis is inferred from the analytic squares.

\paragraph{R-03. Central quotient, disconnected stabilizer, and rank.}
The tensor morphism has the central subgroup $t\mapsto(tI,t^{-1}I)$ in its kernel, a statement of group schemes over $\mathbb Z$. A compact block-Levi branching rule is therefore a different restriction functor. The transpose component of the determinant stabilizer is additional group-scheme data beyond the connected tensor subgroup. The $S_4$ counterexamples compare distinct invariant functors and survive as finite character computations. D-S3 places these groups and actions; cote~133, scan~019, lines~45--56, is the corresponding warning that augmentation loses equivariant information. The negative Weingarten coefficient is a complex integration fact, not a negative rank of a locally free sheaf.

\paragraph{R-04. Plethysm as an operation on representation $K_0$.}
The cycle-index formula expresses $\sigma^d(\sigma^n[E])$ using the Adams operations $\psi^j$. Its displayed rational coefficients are interpreted in $K_0\otimes\mathbb Q$ or the universal torsion-free character ring; the resulting symmetric-power operation is integral. D-S6, V~7.1 and V~7.5, makes their algebraic status precise. They are ring operations; $g\mapsto g^j$ need not be a group homomorphism, and $\psi^j$ need not preserve the cone of honest representations. The $U(2)$ mixed-moment counterexample therefore translates to the statement that separate character distributions do not determine the joint operation. The wreath-subgroup invariants give the characteristic-zero tensor-functor realization. Cote~133, scan~028, lines~12--17 and~26--35, places tensor powers and induction products, with each $\rho_i$ in its own homogeneous degree.

\paragraph{R-05. Image modules, not merely stabilizer invariants.}
The intrinsic object is the image of the degree-$d$ orbit-coefficient map. Its dual is the span of the orbit of $f^{\otimes d}$. Over $\mathbb C$, complete reducibility identifies the multiplicity in each Schur summand with the rank of its multiplicity matrix. The Haar covariance and its reduced Gram matrices have exactly those images; their principal exterior coefficients are squared Pl\"ucker norms. This transports the entire rank proof to determinantal rank strata and Grassmannians in D-FGA. The Hermitian form is an auxiliary real structure and is not asserted to define a motive. Cote~133, scan~028's question about a direct summand is answered here only after passing to characteristic zero.

\paragraph{R-06. Integral matrices and special fibers.}
The polynomial orbit map is defined over $\mathbb Z$ after the specified symmetrizer scalings. Its finite coefficient matrix is an integral map of free modules. The Gaussian covariance is another complex Gram form on its image, while the integral character-sum matrix computes the required Schur-isotypic rank after scalar extension. The Schur projectors require the stated factorial denominators. G-04 specifies the open localization on which the ranks and kernels specialize correctly. D-FGA records the associated rank and Quot strata; D-S6 records the tensor and Schur operations. This translation preserves the exponential matrix sizes, so changing language does not remove the computational gate.

\paragraph{R-07. Binary restriction and highest weights.}
The determinant-pencil substitution is a morphism of coefficient schemes, defined after clearing the displayed binomial denominators. The sections $Q_j$ are highest-weight polynomial functions, and independence is proved by their restrictions to the chosen diagonal-pencil chart. Over $\mathbb C$ every binary form is a product of linear forms, so the binary coefficient subscheme lies in the determinant orbit closure. Highest-weight sections of length at most two therefore cannot belong to its ideal. This is the same proof expressed as ideal restriction in D-FGA and weights in D-S3. Cote~133, scan~028 supplies the polynomial-functor viewpoint; it does not supply a new multiplicity formula by rank alone.

\paragraph{R-08. The quotient class in equivariant $K_0$.}
For the polynomial ring $S=\operatorname{Sym}(W)$ and orbit quotient $A=S/I$, an equivariant free resolution yields $[A]=\sum_i(-1)^i[S\otimes T_i]$ in the appropriate graded Grothendieck group. Symmetric powers produce the Molien denominator and the alternating resolution produces the numerator. Exactness, not pointwise sign, proves the formula. D-S6, V~2.2(b), and D-FGA place the additivity and coherent quotient. Cote~133, scan~019, lines~9--20, contains the exterior/symmetric generating-series relation. Ignoring the numerator forgets the defining ideal and is not an operation allowed by $K_0$ additivity.

\paragraph{R-09. A four-factor invariant and a finite-rank cutoff.}
The squared rectangular multiplicity sum is the determinant-isotypic component of $\operatorname{Sym}(E_1\otimes E_2\otimes E_3\otimes E_4)$ in degree $nd$, obtained by composing two tensor restriction functors. The $U(N)$ cutoff means that Schur functors with more than $N$ rows vanish on a rank-$N$ bundle. The stable ratios denominator is therefore an identity of symmetric operations on four tensor factors. D-S3 and D-S6 carry the proof without an arithmetic purity hypothesis. Cote~133, scan~028, is the corresponding tensor-power locus. Nonnegative character norms are exact, but no effective algorithm for the remaining coefficient extraction has been added by the translation.

\paragraph{R-10. The quadratic case and the component group.}
Over $\mathbb C$, the $2\times2$ determinant is a nondegenerate quadratic form in four variables. Its orbit closure is the whole quadratic-form coefficient space. The connected group is the tensor image of $\mathrm{SL}_2\times\mathrm{SL}_2$; the full quadratic stabilizer includes the other orthogonal component. Accordingly the all-odd four-row part present in the ordinary rectangular Kronecker coefficient is removed from the symmetric coefficient and the coordinate ring. D-S3 records the group and component distinction, while D-S6 identifies the even-part Schur decomposition. Cote~133, scan~019's augmentation does not retain this distinction.

\paragraph{R-11. Low-degree surfaces and restricted ideals.}
The linear and quadratic determinant constructions are coefficient-scheme morphisms. The cubic statement uses a dense open set of smooth cubic surfaces and Beauville's determinantal representation theorem; closedness extends the containment. Block padding is a morphism over the coefficient ring, and the highest-weight ideal inclusion is the same D-S3/D-FGA argument as G-08. It excludes four-row multiplicity obstructions for padded forms of degree at most three without asserting a global orbit containment for arbitrary many-variable cubics. Cote~133, scan~028 provides the Schur-functor interpretation of the resulting coordinate modules.

\paragraph{R-12. A Jacobian certificate and its characteristic-zero consequence.}
The permanent-pencil coefficient map is a morphism of affine schemes over $\mathbb Z$. A maximal Jacobian minor nonzero modulo the stated prime is a nonzero integer and therefore nonzero over $\mathbb Q$. In characteristic zero the differential-rank criterion gives dominance, so the four-variable projection of the permanent orbit closure fills the coefficient space in the certified sizes. D-FGA then translates the hypersurface quotient sequence of G-07 into the exact four-row multiplicity difference. The finite-field calculation is a certificate for an integer minor; it is not an appeal to equidistribution or an assertion of semisimplicity modulo that prime. The tensor-power language is cote~133, scan~028; the rank/character distinction remains scan~019.


\paragraph{R-13. A finite family of integral rank certificates and a torus quotient.}
For each stated size the Jacobian is a map of free modules on the affine pencil-parameter scheme. Its displayed minor defines a principal open rank locus over $\mathbb Z[1/101]$ and over $\mathbb Q$. The independent permanent identities check that these are the derivatives of the same morphism, rather than arbitrary matrices. For larger sizes the row-column torus acts in the fibers with the computed one-dimensional stabilizer. D-S3, VI$_B$~4.1 and the cited EGA dimension result, together with the elementary fiber-dimension calculation, give the upper bound. Highest-weight ideal restriction in D-FGA then yields the stated finite-range multiplicities. This extends the rank and Schur-functor dictionary of cote~133, scan~028; it does not turn a modular certificate into a characteristic-free semisimple decomposition.

\paragraph{R-14. The projective multiplication image and Hilbert growth.}
The map sending a line and a degree-$m$ form to their padded product is a morphism from $\mathbb P^3\times\mathbb P(\mathcal P_{m,4})$. Properness makes its image closed, and the unique linear factor of a general product makes the morphism generically injective. The resulting affine cone has dimension $c_m+3$. Its homogeneous ideal and Hilbert function are the D-FGA objects of EGA~II~2.4--2.5 and FGA Expos\'e~221, Theorems~3.1--3.2. The dimension comparison forces a larger Hilbert function in large degree, and D-S3 identifies the corresponding highest-weight multiplicities. The same proof therefore transports exactly; no effective degree threshold or new purity statement is inserted by this reformulation.

\paragraph{R-15. A normal direction to a family of pencil images.}
The cyclic pencil is a point of the affine parameter scheme, and $W_0$ is the image of its tangent map. Restriction to the explicitly defined lines $\Lambda_i\subset\mathbb P^3$ and the independent interval sections computes a lower bound for that image. The residue calculation proves that the section $T$ gives a nonzero class in the quotient by $W_0$. The two-chord family makes this class lie in the tangent image away from the special parameter, while an augmented determinant witnesses persistence of independence on a principal open. These are explicit maps of coefficient modules and open rank loci in D-FGA, defined over an integer coefficient ring after choosing the stated integral parameters. The generic characteristic-zero dimension and Hilbert argument then gives the four-row inequality for every equal size. D-S3 and D-S6 identify its weight and Schur class. Cote~133, scan~028 supplies the tensor-module setting; scan~019 explains why this is a multiplicity assertion rather than a numerical CUE dimension alone.

\subsection{Isolation as a torus action and a problem of coefficient modules}

\paragraph{I-01. Independent variable coordinates.}
The vertex-variable path polynomial is a section of the structure sheaf of an affine toric coordinate space. Its path monomials are distinct characters of its split torus. A Hoza weight is a cocharacter, and uniqueness of the minimum produces a one-dimensional initial character with a nonzero coefficient. D-S3 identifies the action and grading; I-10 supplies its flat hypersurface family over a field. Cote~133, scan~019, is only a finite-set analogy here, not a claim that path existence already computes a representation multiplicity.

\paragraph{I-02. Descent through a monomial quotient.}
The monomial map is a map of affine toric schemes whose character-lattice map is the integer matrix $A$. A cocharacter is induced by an integral weight on the full variable torus exactly when its degree vector lies in the image of $A^{\mathsf T}$. Homogeneity of the binomial kernel gives the grading on the image algebra; extending it to the full variable torus additionally requires the integral image-lattice condition in I-02. For example $z\mapsto x^2$ has zero kernel, but weight one on $z$ is not induced by an integer weight on $x$. The diamond is an explicit failure even of kernel homogeneity. This is D-S3, I~4.7.3, applied to the quotient algebra. Cote~133, scan~019 supplies no rule permitting independent path gradings to be identified with the quotient grading.

\paragraph{I-03. A complete coefficient flag.}
Filter the coefficient module by the span of monomials below each cocharacter weight. Basis isolation says that at each growth degree precisely one indexed coefficient survives, and that these surviving degrees are distinct. The dual formulation quantifies over every nonzero quotient functional. This is a statement about filtered locally free modules over a field and their graded pieces, not just the presence of one initial path. D-S3 gives the associated grading, and D-FGA places the open minor conditions on the relevant flag and Grassmann schemes. Cote~133, scan~028's direct-summand question motivates the distinction between a filtration and a chosen splitting.

\paragraph{I-04. The Pl\"ucker weight polytope.}
Take the coefficient row plane in a Grassmannian and its nonzero Pl\"ucker coordinates. Their torus weights are the sums of the corresponding monomial exponents. The vertex-fiber assertion is the symmetric-exchange proof stated in I-04; the probabilistic estimate counts walls crossed in one coordinate. The resulting open cones are normal cones refined by basis-weight inequalities. D-S3 identifies their cocharacters and D-FGA, Expos\'e~221, Proposition~3.8 places the Grassmannian. No finite-field purity theorem chooses an integer point in all these cones automatically.

\paragraph{I-05. Real Gram positivity and algebraic minors.}
With two independent dual coefficient matrices, Cauchy--Binet is an integral identity of sections of a determinant line. On the real locus where one is the conjugate transpose of the other, it becomes a sum of nonnegative squared Pl\"ucker coordinates. The counterexample shows that this determinant line does not retain the whole coefficient-image flag after colliding weights. D-FGA supplies the rank-stratum language, and D-S6 supplies the determinant as top exterior power. Cote~133, scans~019 and~028, respectively mark the distinction between numerical positivity and actual tensor-module information.

\paragraph{I-06. Secant rank of an evaluation functional.}
An evaluation at a point is a rank-one decomposable functional on the bounded-degree polynomial space. The moment matrix of the proposed nonlinear exponent functional has tensor-product rank. Expressing it as a fixed sum of evaluations would place it in a secant locus of bounded rank; the computed determinant excludes that locus below the stated exponential size. This is an algebraic rank statement over $K(t)$, hence belongs to the determinantal geometry in D-FGA. D-S3 also explains why a nonlinear exponent function is not a character-lattice homomorphism. No motive is needed for this counterexample.

\paragraph{I-07. Taking invariants changes the inserted data.}
The invariant functor on complex representations is realized by the Reynolds projector. An independent twirl sends each state insertion to its invariant scalar component. The two-insertion calculation resolves the commutant into the identity and transposition operations and exhibits a negative coefficient. Thus invariants are extracted in the representation category of D-S3/D-S6, with the full tensor operation retained. Cote~133, scan~028, describes why the symmetric-group action on tensor powers matters; scan~019 warns against replacing that action by dimension.

\paragraph{I-08. A finite quotient morphism, conditional on a hitting set.}
The trace-generating functions are regular invariant functions on the affine scheme of matrix tuples. A hitting set separating their coefficient polynomials gives a separating family in the quotient coordinate algebra; the cited invariant-theory theorem makes that algebra finite over the generated subalgebra. This is a finite morphism of affine schemes, not automatically a choice of algebraically independent parameters of minimal cardinality. D-FGA gives the general morphism and parameter-space setting. The unresolved explicitness condition is exactly the one retained in I-08 and GCT~V; the dictionary does not discharge it.

\paragraph{I-09. A tensor contraction computing a Hermitian norm.}
The matrix product computes a regular coefficient map. Tensoring with its conjugate and contracting equal monomial indices gives its squared norm; the torus integral extracts the zero character in that tensor product. D-S3 explains the character extraction. Over an integral algebraic base one must use independent dual variables, not treat complex conjugation as a regular morphism. The metric realization lives over $\mathbb R$ or $\mathbb C$. It is the same coefficient module as before, and access to its presenting matrices remains a white-box assumption.

\paragraph{I-10. Rees modules and specialization.}
The cocharacter gives an increasing filtration on coefficient spans. Its Rees module has the associated graded as its special fiber and the original module on the punctured line. The explicit hypersurface homogenization in I-10 is flat over a field by the torsion-free PID argument in its proof; general ideal degenerations require the entire weight-homogenized ideal. D-S3 identifies the grading and D-FGA identifies the flat family and its Hilbert/Quot points. The associated graded preserves total filtered rank, whereas adding colliding coefficients in a substituted polynomial can reduce their span, as I-05 proves. These are different maps.

\paragraph{I-11. Exterior powers do not remove equations.}
The exterior dimension and graph-of-linear-map count describe the sizes of coordinate and Grassmannian models. The two-term zero determinant is already an equation of the rank-deficient locus. Assigning different path weights to its cancelling terms changes that equation. D-S6 gives exterior powers and D-FGA gives the Grassmannian, but neither turns a signed determinant expansion into unsigned reachability. Cote~133, scan~028's direct-summand margin remains a relevant distinction between a functorial module and its termwise expansion.

\paragraph{I-12. Support is the nonvanishing locus of a regular function.}
A specified coefficient of an integer ROABP is a regular polynomial in its edge parameters. Its nonvanishing open subset is determined by a signed path sum, whereas ordinary reachability detects whether a parameterizing path set is empty. The GapL reduction is a computation of that regular function; the conditional complexity-class inclusion follows only under the specified logarithmic-space reduction. D-FGA locates its zero scheme. Cote~133, scan~019 supplies the nearby conceptual distinction between the cardinality of a finite set and a signed class; it states no complexity theorem.

\paragraph{I-13. Evaluation quadrature and the Segre variety.}
Coefficient tensors of width-one products form the affine cone over a Segre variety. Equality of the quadratic forms on all these tensors, with their complex conjugates, forces the moment-matrix identity proved in I-13. Its rank then bounds the number of evaluation vectors. This restates the entire proof as an assertion about a Segre cone, its dual evaluation vectors, and determinantal rank loci in D-FGA. The Hermitian positivity is not an integral $\lambda$-ring assertion. The much smaller hitting set in the same finding concerns only kernel avoidance, not equality of quadratic forms.


\paragraph{I-14. A Grassmannian open condition for several substitutions.}
The row plane of the coefficient matrix is a point of $\operatorname{Gr}(r,F^s)$. A fixed family of cocharacters gives the integral linear map $P_{\mathcal W}^{\mathsf T}$. The desired locus is the open set where its restriction to the tautological rank-$r$ bundle is injective. Its complementary equations are the displayed linear forms in Pl\"ucker coordinates. D-FGA, Expos\'e~221, Proposition~3.8 is the Grassmannian framework, and D-S3 identifies the cocharacters. The example proves that this open set depends on coordinate values and can be larger than the union of the individual basis-isolation loci.

\paragraph{I-15. A finite character-lattice obstruction.}
The bounded seed decoder produces a subspace of cocharacters of the split torus. Pigeonhole counting gives two distinct characters whose difference annihilates every cocharacter in that subspace. Their binomial is therefore a nonzero regular function in the kernel of every pullback. This is precisely the character/cocharacter pairing of D-S3, I~4.7.3, with the finite integer bounds retained. The extension to bounded-degree Boolean decoders is a dimension bound for their cocharacter span. It does not forbid more general maps whose spans are larger.

\paragraph{I-16. A finite coefficient-state scheme and base-field points.}
Over $\mathbb F_q$ the coefficient row space has $q^w$ rational points, and each letter acts on it by a linear morphism. The state automaton records these morphisms exactly. The separate evaluation lower bound counts the rational points avoiding products of prescribed coordinate divisors. Both constructions live in affine schemes over the finite field; D-S3 gives the coefficient representation, while D-FGA distinguishes the finite set of rational points from the whole scheme or its extension-field points. Purity does not change the explicit counting inequality.

\paragraph{I-17. A trivialization on an invertibility open set.}
The nonzero determinant polynomials define an open set on which the constant transition maps are bundle isomorphisms. The displayed gauge change trivializes those constant maps. The nested coefficient row subspaces can then grow at most $w-1$ times, and their chosen sections have the stated support. This is a concrete filtered-module argument in D-S3 and a rank-stratum argument in D-FGA. The known AGKS invertible-factor result is cited separately; the proof establishes the particular permutation-invariant sampler stated here.

\paragraph{I-18. A unit minor over a discrete valuation ring.}
Over $F[[\varepsilon]]$, a maximal minor defines an open rank locus precisely when it is a unit. A nonzero nonunit meets the generic fiber and misses the special fiber. The explicit matrix family separates these two cases, and the ideal of sampled scalar values shows why saturation in the fraction field does not recover an original evaluation. This is the elementary local-module form of the determinantal loci in D-FGA. The distinction from a flat Rees degeneration is essential: flatness of an ambient family does not imply surjectivity of this particular compressed module map.

\paragraph{I-19. Rank-factorization charts and the coefficient field.}
On the principal open where a chosen minor is invertible, a rank-$r$ bundle map factors through a rank-$r$ bundle. The adjugate formula clears its denominator and is an integral polynomial identity. Passing to the function field of the cut-variable scheme makes the intervening blocks disjoint; it does not remove the cut variables from the original base. Multilinearity proves coefficient-span preservation under independent local substitutions. D-FGA supplies these rank charts and D-S3 the tensor contractions. The unresolved descent from the cut-variable function field is explicitly retained.

\paragraph{I-20. Bounded losses in a sequence of coefficient subspaces.}
The product of a nonzero coefficient and the selected pivot minors defines a principal open on the cut-coordinate affine scheme. The finite grid meets this open by the degree bound. After specializing there, the coefficient subspaces have total dimension loss at most $\kappa$; hence their total growth is at most $w-1+\kappa$. This is an exact rank and filtration argument on the same D-FGA charts as I-19, followed by the D-S3 coordinate grading. It yields the stated sampler with its parameter dependence, without asserting that arbitrary singular sequences have bounded total loss.

\subsection{Purity, arithmetic sampling, and the limits of numerical shadows}

\paragraph{P-01. A Gram form on a finite space of sections.}
The coefficient support spans a vector space of functions. Their pairwise correlations are entries of a Gram matrix; its spectral lower bound is the elementary estimate in P-01. In geometric applications these functions may be sections or Frobenius traces, but no sheaf is manufactured by naming the Gram matrix. D-ET applies only after an actual trace-function realization exists. The distinction from rank augmentation in cote~133, scan~019, is that all cross terms have been retained and bounded.

\paragraph{P-02. A Kummer-character construction.}
The signs arise from the rank-one quadratic Kummer sheaf pulled back by $t\mapsto t+i$ on the line with its zeros removed. Products over distinct indices remain governed by the explicitly squarefree polynomial, so the curve character estimate controls the off-diagonal trace sums. The altered values at the excluded points contribute the stated elementary error. D-ET, together with the exact Alon--Goldreich--H\aa stad--Peralta locator, carries the proof. This yields the stated sparse multilinear hitting set; it does not create a polynomial bound on arbitrary coefficient supports.

\paragraph{P-03. The kernel of restriction.}
The sign-cube counterexample is a nonzero section whose restriction vanishes on the selected finite subscheme; the common-affine repair is annihilated by the explicit Vandermonde section. These are regular-function identities and survive independently of improved point-counting estimates. They belong to the affine-scheme language of D-FGA. Cote~133, scan~019 is an analogy for information lost by passing to a smaller numerical image, not a source of these new counterexamples.

\paragraph{P-04. Pullback to a curve and its zero divisor.}
For the curve morphism specified by the rational functions $r_i$, the section $f(r_1,\ldots,r_n)$ has polar divisor bounded by the displayed $D$. If it is nonzero, its zeros have total degree at most $\deg D$. The trace formula and curve RH then provide enough rational points. D-ET controls the count; the pullback's nonvanishing is a separate scheme-theoretic condition. This proof uses curves and sheaves actually specified in the statement, so a motivic interpretation is available for their cohomology, not for an unspecified ROABP.

\paragraph{P-05. A curve family can lie in a small polynomial kernel.}
Dimension comparison in the spaces $H^0(C,\mathcal O_C(aD))$ produces an algebraic relation between the two chosen functions. Multiplying the relations for the finite curve family gives its annihilating section and the explicit bivariate ROABP. D-FGA places the image and its ideal; D-ET cannot remove a function from that ideal by increasing the field size. The relation is proved directly with its degree and width bounds and is not a lower bound against every geometric sampler.

\paragraph{P-06. Purity with a complete conductor budget.}
The relevant class is $[\mathcal F_i\otimes\mathcal F_j^\vee]$ in the Grothendieck group of constructible Weil sheaves. The compactly supported trace gives its Gram entry. Geometric irreducibility and inequivalence determine $H_c^2$, affine vanishing removes $H_c^0$, and Weil~II bounds $H_c^1$. SGA~5 supplies the displayed Swan-conductor expression for its dimension. This is exactly D-ET, with every rank and conductor retained. A claim of polynomial explicitness additionally requires that these data and the enumeration of the points be polynomially bounded. Cote~133, scan~028 supplies a tensor-category analogy, but not the missing sheaf realization.

\paragraph{P-07. Sparse cocharacters over the integer lattice.}
The weight family is an explicit finite set in $\operatorname{Hom}(\mathbb G_m,\mathbb G_m^n)=\mathbb Z^n$. The collision polynomial over $\mathbb F_p$ is nonzero by its degree and coefficient bounds, so a seed outside the union of its root sets makes the support characters distinct. D-S3 translates the construction exactly; only the elementary finite-field root bound is used. Thus the result is independent of an additional purity or RH hypothesis.

\paragraph{P-08. Arithmetic covers with a specified nonvanishing encoding.}
The common number field determines a finite \emph{\'etale} cover of $\operatorname{Spec}\mathbb Z$ after inverting its discriminant. Frobenius conjugacy classes are the arithmetic data on that cover. Effective Chebotarev produces a small unramified prime only after the stated common map from primes to evaluations and its nonvanishing property are supplied. D-ET provides the parallel finite-field trace perspective; it is not a proof of that arithmetic encoding. The abelian case uses the stipulated Dirichlet GRH, while the general stated Hecke extension remains conditional on Masullo. No claim about arbitrary automorphic $\mathrm{GL}_n$ GRH is made.


\paragraph{P-09. Positive Hermitian forms on a Segre cone.}
The width-one coefficient tensors are decomposable sections, forming a Segre cone. Applying the upper sampled-energy bound to the conjugate evaluation tensor bounds each weighted evaluation norm, while applying the lower bound to the monomial sections bounds their total trace. This is the complete rank-one tensor proof of the cardinality inequality. D-S6 identifies tensor products and D-FGA the Segre coefficient scheme. The inequalities themselves use a complex Hermitian realization and are not inequalities in integral $K_0$.

\paragraph{P-10. Projective incidence and avoidance of a determinantal cone.}
The tensor-train class is the closed cone defined by the stated flattening minors. Its projectivization and evaluation hyperplane sections give a projective incidence correspondence. Proper projection makes the bad tuples a closed set; successive hyperplane sections show this set is proper at the stated cardinality. For unknown order one uses the finite union of the permuted cones. This is D-FGA's projective and determinantal geometry with D-S3's coefficient grading. Compactness of the complex unit sphere gives a positive Hermitian lower bound only after kernel avoidance is known. Neither the projective argument nor the $K_0$ class bounds the integer heights of an explicit good tuple.

\paragraph{P-11. Pullback along an explicit arithmetic parameterization.}
The presentation $\alpha_i=A_i(\theta)/b$ specifies a morphism from an affine parameter scheme to the variable space after inverting $b$. Pulling back a polynomial to $G_f(U,T)$ and using the stated degree grid gives the unconditional forward construction. Conversely Lagrange interpolation and a cyclotomic point give a finite \emph{\'etale} arithmetic realization after inverting the discriminant; minimal-polynomial degree prevents a nonzero rational pullback from vanishing there. D-FGA describes the morphism and its kernel, and D-ET places the finite cover. The proof does not require a Frobenius estimate: the missing hypothesis is the explicit common nonvanishing map itself, rather than a further trace bound.

\paragraph{P-Z01 and P-Z02. Realizations of zero statistics.}
These are analytic consequences of the precisely stated RH predicates and supported explicit-formula theorems. Their expression by unitary integrals does not make the zeta zeros the eigenvalues of a constructed geometric Frobenius. D-ET explains the function-field comparison only in the families with an established sheaf and monodromy group. There is no cote~133 scan containing a classical-RH-to-GCT implication; scans~019 and~028 supply representation operations, not that missing implication.

\subsection{The integral controls G-01--G-12}

\paragraph{G-01. Burnside classes, representation classes, and rank.}
The Burnside-ring counterexample is an equivariant finite-set computation. Linearization sends a $G$-set to a permutation module, but the two-subset operation does not become exterior power: the $C_2$ sign in the proof witnesses this failure. In D-S6 the correct relation is $x^2=2\lambda^2(x)+\psi^2(x)$, and only augmentation replaces $\psi^2(x)$ by the rank. This is the precise mathematical correction to cote~133, scan~019, line~53 and the supplied study's equivariant reading.

\paragraph{G-02. The universal terminal variable.}
Specializing one factor to the tensor unit in the universal product identity forces the top elementary symmetric function to remain present. SGA~6, V~2.4, equation~(2.4.2), explicitly includes that terminal variable. The proof in G-02 is its necessity in the universal character ring, so the correction to scan~006, lines~25--26, is mathematical and independent of paleography.

\paragraph{G-03. Equal finite filtration data, unequal representations.}
The two sums of line characters are honest locally free representations of $\mathbb G_{m,\mathbb Z}$, whose difference is the $(D+1)$st power of an augmentation-zero element. D-S6, V~3.10, places the difference in the required filtration step. The associated-bundle construction transports this identity to any scheme carrying the line bundle; it does not force the bundles to remain distinct on every base. Thus finite Chern or $\gamma$ data cannot recover all the character coefficients needed in scans~019 and~028.

\paragraph{G-04. The integral image and its good rank locus.}
The schematic orbit closure is defined over $\mathbb Z$ by a kernel, and its image algebra is torsion-free. In a fixed degree a nonzero maximal minor gives an explicit principal open on which the coefficient matrix is equivalent to a split rank matrix. This is a direct proof of base change for the displayed module, with D-FGA providing the image, Grassmannian and Quot framework. At the excluded primes the flat integral model's fiber may differ from the orbit closure of the reduced point. This is why integral $K_0$ identities from scan~028 do not imply unchanged characteristic-zero multiplicities in all fibers.

\paragraph{G-05. Plane curves and the three-row sector.}
The nonempty complement of the theta divisor in $\operatorname{Pic}^{g-1}(C)$ supplies Beauville's line bundle, hence a linear determinantal equation. Density of smooth curves gives the scheme-theoretic containment needed for every ternary coefficient section. Highest-weight sections of length at most three belong to that coefficient subalgebra, so its determinant ideal has none. D-FGA records the ideal restriction and parameter spaces; D-S3 records the weights. The numerical positive count is not extracted from Chern classes of the curve alone.

\paragraph{G-06. Stabilizer dimensions and four-row equations.}
The polynomial matrix-pencil map is invariant under the explicitly stated subgroup scheme $H$ of $\mathrm{GL}_n^2$. Its generic stabilizer is the scalar subgroup by the concrete matrix calculation. D-S3, VI$_B$~4.1 and its cited EGA theorem justify the dimension behavior, and the resulting proper image has a nonzero ideal. D-FGA and highest-weight restriction give an equation of exactly four rows, using G-05 to exclude shorter weights. This supplies no nonvanishing on a particular padded form until that extra statement is proved.

\paragraph{G-07. A Cartier divisor and a determinant twist.}
The reduced determinantal-quartic locus is a prime hypersurface. Its equation spans the character $\det^D$, so its ideal is the line-character twist of the ambient graded algebra shifted by $D$. The exact sequence in G-07 is the equivariant $K_0$ statement of a principal Cartier divisor. D-S6 gives additivity, D-FGA gives the projective hypersurface and quotient, and D-S3 identifies the character. It produces the shifted plethysm subtraction exactly; the corrected reduced degree is a literature issue, not a new transcription of cote~133.

\paragraph{G-08. A universal morphism of coefficient schemes.}
The two monomial trees and their cofactor determinant give a matrix over the universal integer coefficient ring. Homogenization and block padding therefore commute with any base change. Over $\mathbb C$ the containment of all low-variable restrictions yields an inclusion of highest-weight ideal spaces, hence the multiplicity inequality. This is D-FGA applied to an explicit scheme morphism, with the weights read through D-S3 and the symmetric powers through D-S6. Its complexity bound is part of the construction, not a consequence of representability alone.

\paragraph{G-09. A local padding coordinate and smaller universal matrices.}
Once the nonzero padding form is one coordinate, dehomogenization reduces the number of affine variables by one. The shared exponent chains are again defined over the universal coefficient ring; their exact vertex count gives the stated finite size. Changing an arbitrary padding form to that coordinate is local on its nonzero coefficient charts. D-FGA carries these local morphisms and their closure statement, and the D-S3 highest-weight argument yields the improved row exclusion. The known asymptotic uniform-representation theorem is cited separately; neither the chart construction nor its GCT consequence requires RH.


\subsection{Further rectangular, rank-profile, and arithmetic consequences}

\paragraph{R-16. Several normal directions on one deformation.}
The cyclic pencil lies on the coefficient affine scheme; disjoint row-pair modifications define one morphism from the affine line into that parameter scheme. The exact permanent identities produce polynomial sections of its differential-image sheaf away from zero. Their special-fiber classes are separated by characters of the finite diagonalizable group of $m$th roots of unity after adjoining $\zeta$ and inverting $m$. Dual Jacobi--Trudi is the Schur operation in D-S6, equivalently the exact cote~133 symmetric/exterior dictionary. A nonzero maximal minor makes the normal directions independent on a principal open. D-FGA and D-S3 therefore carry the dimension and padding proof; the characteristic-zero conclusion does not assert uniform good reduction at every prime.

\paragraph{G-10. Projective degree controls the first forced Hilbert gap.}
The pencil map induces a rational morphism of projective coefficient schemes. Pullback of general hyperplanes and the degree of their intersections bound the image degree, including components of the base locus. Chardin's bound then compares the graded dimensions with a polynomial subalgebra supplied by linear Noether normalization. These are D-FGA's projective, graded-module and Hilbert frameworks, with associated sheaves as in EGA~II, Proposition~2.5.2. The explicit bound is geometric effectivity; no small-size representation of every graded module is inferred.

\paragraph{G-11. Rectangular classes in a reductive quotient.}
Restriction from $\mathrm{GL}_4$ to $\mathrm{SL}_4$ sends precisely the determinant characters among polynomial irreducibles to the tensor unit. In D-S3 and D-S6, the degree-$d$ invariant module is therefore the rectangular multiplicity space. Explicit moment-map witnesses establish a nonempty stable open locus over $\mathbb C$; its quotient has dimension reduced by $15$, by the fiber-dimension statement already located in D-S3. The invariant rings are quotients because the group is linearly reductive in characteristic zero. The explicit Veronese is a standard graded coordinate algebra, so D-FGA applies again to its Proj, degree and Hilbert function. The analytic compact-group witness proves the stated characteristic-zero fact; it is not asserted to define a positive integral motive or to make reductive invariants exact over $\mathbb Z$.

\paragraph{G-12. A cocharacter degeneration and its invariant ideal.}
The homomorphism $t\mapsto\operatorname{diag}(t^3,t^{-1},t^{-1},t^{-1})$ is an integral cocharacter of $\mathrm{SL}_{4,\mathbb Z}$, as in D-S3. Its weight filtration extends the padded form over the affine line whenever $3(n-m)\geq m$. Strict inequality gives the zero special fiber; equality gives a ternary residual form with its padding factor. The ternary determinantal containment of G-05 then implies containment of invariant ideals in characteristic zero. In D-FGA this is the specialization of a morphism and the inclusion of its ideal of functions; in D-S6 it selects the determinant-character summands. It excludes precisely the specified rectangular sector.

\paragraph{I-21. Rank strata over local Artin bases.}
Every interval coefficient product is a map of free modules. The proof works over finite local Artin rings and requires a unit maximal minor, so the rank-two image and its complement split on the corresponding determinantal chart. Projection growth and its last-zero reset control the necessary coefficient supports. The two selector stages pull these finitely many rank conditions back to one nonempty open parameter locus. D-FGA supplies the determinantal and Grassmannian language, while D-S3 translates the coefficient supports to character gradings. This gives the explicit unknown-order width-three morphism; it does not assume that an arbitrary rank stratum over a nonreduced base splits without a unit minor.

\paragraph{P-12. A finite flat arithmetic evaluation and its trace dual.}
The algebra $\mathbb Z[T]/(T^N-2)$ is finite free of rank $N$. After inverting its discriminant it is finite \emph{\'etale}; over $\mathbb Q$ it is a field by Eisenstein. The integral coordinate polynomials $A_i$ define the evaluation morphism. Their bounded-degree pullback does not annihilate any nonzero member of the specified ROABP class, by I-21 and the explicit monomial encoding. The dual basis for the trace pairing extracts the transformed rational coefficient rows, which span the original coefficient-row space. This is an actual span-preserving arithmetic realization, not an injectivity assertion for differences of arbitrary members of the nonlinear width-three class. D-FGA describes the affine finite morphism and its kernel, D-S3 the coefficient lattice, and D-ET the finite \emph{\'etale} cover. Frobenius distribution and additional purity bounds are unnecessary for this construction.


\paragraph{R-17. Identical tangent-image modules and distinct isotypic sectors.}
At the cyclic pencil, the determinant and permanent differential maps have the same image up to signs of their columns. Their rank equality, together with the group-scheme fiber bound, gives the exact determinant-image dimension. D-FGA and D-S3 transport this computation. The two finite padded pairs then compare Hilbert polynomials, while the cocharacter specialization of G-12 removes the determinant-character summands. Their remaining positive differences are necessarily nonrectangular Schur summands in D-S6.

\paragraph{R-18. Two compatible normal families on the same scheme curve.}
Both local permanent identities hold on the same one-parameter deformation. The character decomposition of the special-fiber quotient now has two independent sections per permitted mode. The four-column coefficient determinant is the Schur operation $s_{(4^b)}$, interpreted in D-S6; its root-of-unity specialization is checked by a Weyl alternant. D-FGA turns their joint nonzero minor into an open rank condition, giving the improved padding range. The simultaneous nature of the deformation is necessary: ranks obtained at unrelated pencils are not added.

\paragraph{G-13. A weighted filtration inside an ordinary projective embedding.}
Generators of the reductive invariant quotient have specified positive degrees. Its cumulative graded module embeds in the ordinary polynomial-degree filtration of their affine coordinate embedding. Homogenizing that embedding gives an integral projective closure, to which the D-FGA Hilbert and degree comparison applies. Algebraically independent homogeneous sections provide the lower filtered dimension. This reduces the finite first-gap bound without asserting equality of the two filtrations or generation by a small Veronese degree.

\paragraph{I-22. Ordered coefficient modules and the radical quotient.}
The actual ROABP coefficients occupy an ordered product subspace, which need not equal the associative algebra generated by all layer coefficients. Passing to the semisimple quotient loses radical-valued sections even when their specified boundary functional is nonzero. In D-FGA these are different linear maps and different kernels; in representation $K_0$ the composition factors alone do not recover the off-diagonal extension data. The explicit triangular examples show the missing information before any trace or purity estimate can be applied.

\paragraph{I-23. Cocharacters preserving a Pl\"ucker subscheme.}
The Grassmann cone is cut out by the Pl\"ucker ideal. A diagonal cocharacter preserves that subscheme precisely when its weights satisfy the proved exchange equalities, which identify the integral torus weights induced from the original rank-$w$ bundle, up to a common character. D-S3 supplies the cocharacter lattice and D-FGA the Grassmannian. The binary classification and code-distance sum then quantify the cost of imposing this exact subscheme-preservation condition on every Hoza seed.

\paragraph{I-24. Constant-rank maps with repeated loss of remembered directions.}
The nested projections and their univariate perturbations define maps of free modules over an affine parameter scheme. Every interval lies on its specified unit-minor rank chart, including over the finite Artin quotients used in I-21. Nevertheless the selected boundary section is one monomial of exponential support. D-FGA's rank stratification records this phenomenon exactly; the ranks alone do not bound the length of a coefficient-support filtration.

\paragraph{I-25. Nested projective images on a Grassmannian.}
Invertible factors act regularly on the Grassmannian of rank-$r$ quotients or subspaces. The irreducible prefix image closures are nested because the constant factor is the identity; every strict inclusion raises dimension. D-FGA's Grassmannian chart gives dimension $r(w-r)$, and its projective geometry identifies scalar vanishing with containment in a hyperplane. Removing equal-image steps therefore gives the stated support existence bound. The invertible-factor hypothesis is what allows this proof to omit the reset behavior in I-24.

\paragraph{R-19. Compatible row and column image sections.}
Three-row motifs retain both supported row and supported column factorizations in one parameter family. Their four normal sections are separated by a six-column character calculation, whose Toeplitz determinant is the Schur operation $s_{(6^b)}$. D-S6 gives that operation, D-S3 the diagonalizable character splitting, and D-FGA the common nonzero determinantal minor. The resulting improved range is a statement about one actual pencil deformation; no addition of ranks from incompatible families is used.


\paragraph{G-14. Truncated coefficient jets and Schur-weight exclusion.}
The restriction modulo the square of the padding coordinate defines a linear
jet projection between affine parameter schemes. At a cyclic ternary pencil,
the differential is surjective: its two blocks are the cofactor-span map and
its multiplication by linear forms. The nonzero differential minor gives a
dominant morphism in characteristic zero (D-FGA). Every weight in the Schur
functor of shape $\lambda$ uses any fixed basis direction at most
$\lambda_1$ times, by its exterior-column construction in D-S6. After padding,
the low-weight part of an equation therefore factors through the first jet.
Its vanishing on the dominant determinant-jet image forces vanishing on all
padded forms. This is containment of the corresponding isotypic ideal
submodules, followed by exact passage to the characteristic-zero quotient;
it is not an assertion that the full padded varieties are contained.

\paragraph{I-26. Determinantal equations of an interpolation image.}
The maps $U$ and $V$ define bundle maps with $UV^T=0$. The determinant of
$U\operatorname{diag}(X)V^T$ is a nonzero section by a maximal-minor chart,
but its pullback to the centered interpolation parameter scheme vanishes
because that pulled-back bundle map factors through a rank-$k$ bundle.
The common center acts trivially. All exterior minors form a finite free
module on which rank-one additive translations act by commuting unipotent
matrices; the proof uses polynomial identities and remains valid in the
stated characteristics. This is the D-S3 additive group action and the
D-FGA determinantal rank locus. The surviving column-matroid bases give the
separate cocharacter nonvanishing certificate of D-S3. No purity assertion
turns the annihilated interpolation pullback into a nonzero section.


\paragraph{G-15. Harmonic covariants make the third-jet morphism dominant.}
The $\mathrm{SL}_2$ representation-ring decomposition of the endomorphism
bundle and the symmetric algebra of its adjoint bundle gives the covariants
$T_l$ in D-S3/D-S6. Their scalar coefficients in the adjugate are determined
by a finite trace pairing and its Jacobi matrix. The strictly positive
squared norm in the proof is an ordinary Hermitian linear-algebra fact,
not a weight-purity assertion. Linear multiples of these cofactor sections
surject onto the quotient by the fourth power of the padding-coordinate
ideal. Thus the differential minor gives dominance of the corresponding
third-jet morphism in D-FGA, while dimension excludes the next full jet.
The D-S6 exterior-column weight bound transfers dominance to containment
of the specified isotypic ideal submodules. This closes the precise
higher-jet gate stated conditionally in G-14.

\paragraph{I-27. A Boolean face in the cocharacter lattice.}
The Pl\"ucker ideal of the D-FGA Grassmann cone supplies integral linear
equations in the diagonal cocharacter lattice of D-S3. Under the actual
base-six code embedding, bounded integer deficits cannot carry between
digits. Its intersection with that lattice condition is consequently a
coordinate face of the seed cube. The induced original-bundle characters
then agree on two untested basis directions. The displayed rank-two
unipotent representation detects this equality as exact vanishing of a
nonzero section. No statement about arbitrary decoders follows.

\paragraph{I-28. A hyperplane dimension drop pays for the final parameter.}
The left endpoint defines a curve through its nonzero constant point of
the Grassmannian. Including it in the nested projective-image chain uses
the same dimension budget as the internal bundle automorphisms. If the
constant right functional vanishes on the final image, that image lies
in a proper hyperplane section and loses at least one dimension; this
pays for retaining the right endpoint parameter. D-FGA therefore gives
the sharper $r(w-r)$ support existence statement with both endpoints
included. Its width-three specialization has eight SV pairs, and the
finite flat trace construction of P-12 uses the corresponding seventeen
parameters when sharpened by this result.


\paragraph{R-20. Second-order tangent sections and the exact eventual rectangle range.}
The permanent coefficient morphism has a family of differentials over
the one-parameter motif scheme. Row and column relations cancel its
constant and first-order normal terms. Dividing their polynomial
sections by the square of the parameter gives regular limiting sections
of the ambient coefficient bundle. The finite certificate proves a
nonzero minor after one residue-field specialization; the explicit
degree bounds and cyclotomic argument lift that nonvanishing to the
required characteristic-zero family. D-FGA records the resulting image
dimension, and the D-S3 Fourier characters separate its translated
motifs. The invariant quotient and Veronese construction of G-11 then
give the specified rectangular Schur classes. This transports the
actual proof, including its finite threshold, rather than treating a
dimension inequality as an unspecified positivity assertion.

\paragraph{R-21. Integral image ranks and fixed-rank character extraction.}
The orbit-coefficient homomorphism of free modules has an integral Gram
matrix. Its positive-rank characteristic coefficient is a nonzero
integer, giving the displayed analytic spectral separation. A bilinear
coefficient kernel presents the same image ranks without enumerating
the full tensor basis. The diagonalizable group of D-S3 splits that
kernel into weight blocks; the D-S6 Weyl-character identity extracts
each Schur multiplicity from their ranks. D-FGA identifies the ranks
with the degreewise coordinate-image modules. The algorithm is
polynomial only when the bundle rank is fixed and degree is measured
numerically. The metric estimate is ordinary linear algebra, not an
unproved arithmetic slope transfer.

\paragraph{I-29. A persistent active subbundle and one passive rank.}
On the preserved rank-profile chart, constant transition maps split
into an identity on a rank-$r$ active subbundle and invertible or zero
maps on a rank-$(w-r)$ complement. At the last zero map, the vanishing
higher minors give a row graph over the active summand. The stable
projection argument bounds the support filtration of its coefficient
image by $r(w-r)+w-1$. The Grassmannian dimension argument for rank
preservation adds $r(w-r)$ interpolation pairs. These are explicit
rank-stratum and filtration calculations in D-FGA and D-S3, valid with
arbitrarily many resets; the polynomial pair budget retains its
separate sampling exponent.

\paragraph{I-30. Endpoint sections extend to local automorphisms.}
A nonzero constant endpoint section has a split coordinate functional.
The displayed rank-one update extends the section to a bundle
automorphism on the principal open where its determinant is nonzero.
The determinant is a unit at the center. Thus all generic and constant
interval ranks are preserved after the endpoints are absorbed in the
word. In D-FGA this is an explicit trivialization on the required rank
chart, and in D-S3 it preserves the original coordinate support
filtration. The two width-three rank stages consequently use two and
four pairs, respectively.

\paragraph{I-31. Irreducible matrix cones and deletion of stable blocks.}
The closure of scalar multiples of a polynomial matrix image is an
irreducible affine cone in D-FGA. The rank-$s$ chart identifies the
image of its last block with the graph of a map over the first $s$
columns. Equality of two consecutive projected cones allows deletion:
an invertible right action gives a closed subcone of the same dimension,
which is equal to the irreducible cone. This geometric equality
preserves every homogeneous minor and composes under multiplication
on disjoint parameter schemes. Rank descent and split rectangular
endpoints yield the stated recurrence in the original rank-$w$ bundle.
No exponential compound bundle is introduced, but the recursive block
factor remains and prevents a uniform polynomial sampling conclusion.

\paragraph{P-13. A finite flat algebra encoding any centered generator.}
After clearing Lagrange denominators, distinct base-$B$ exponents embed
the bounded parameter monomials in a single polynomial variable. The
Eisenstein quotient by $T^{B^{2K+1}}-2$ is finite free over
$\mathbb Z$; its generic fiber is a separable field and the displayed
trace-dual basis recovers every coefficient coordinate. The
annihilator argument identifies the resulting image span with the
original coefficient-image module. D-FGA gives the finite algebra and
module interpretation, and D-S3 gives the tensor-coordinate transport.
This is the ordinary algebra trace, distinct from the Frobenius traces
of D-ET, and does not acquire a purity or RH hypothesis by analogy.


\paragraph{G-16. An auxiliary dual line selects a nonrectangular Schur class.}
Form the Segre--Veronese cone of the degree-$n$ coefficient scheme with
the $a$th Veronese of the dual rank-four bundle. Its degree-$d$
coordinate module is exactly the original degree-$d$ module tensored
with $\operatorname{Sym}^{ad}V$. The D-S3 invariant functor and the
D-S6 Schur decomposition select precisely
$(ad+md/4,md/4,md/4,md/4)$. The auxiliary dual line cancels the padding
weight of the explicit Fermat point, so the enhanced cone has a stable
point for every positive padding exponent. D-FGA and the reductive
quotient locators in G-11 give the quotient dimensions; the separate
action and parameter degrees give the effective Hilbert bounds.
Thus the same dimension gap determines an actual nonrectangular
highest weight and a degree progression. The proof does not infer a
highest-weight label from a total dimension inequality alone.


\paragraph{G-17. Partial flags classify the stretched directions.}
For a partition direction $\nu\vdash ns$, subtract its last row to
obtain the polynomial highest weight $\eta$. The orbit of its extremal
line is the partial flag scheme embedded by the exterior-power Schur
construction. Its degree-$d$ coordinate module is the Cartan component
$S_{d\eta}V$, as follows directly by dualizing the orbit span of the
extremal tensor power. Pairing this cone with the $s$th Veronese of the
coefficient scheme makes the D-S3 invariant functor extract exactly
the multiplicity at $d\nu$, in original coordinate degree $sd$.
Positive Fermat coefficients balance the moment map precisely in the
strict region $\nu_4>0$, $\nu_1>(n-m)s$. D-FGA gives the quotient
dimension comparison and its effective Hilbert bounds. The two
excluded faces follow from the earlier ideal containments, so this
is an exact stretched-direction criterion, with its dimension
hypothesis retained. D-S6 records every exterior, symmetric and Schur
operation used in the construction.

\paragraph{I-32. A common-center rank chart with exponential deletion cost.}
Each singular factor is lifted to two univariate bundle maps whose
product depends only on the difference of their coordinates. On the
diagonal center those differences vanish, and every partial interval
is a product of nested projections between invertible boundary maps.
Thus all interval ranks are already correct on that D-FGA chart.
Nevertheless a homogeneous matrix-coordinate section is the product
of all paired differences. A coordinate subspace meets its nonvanishing
locus only if it retains one coordinate from every pair. This gives
the exact exponential threshold for preserving the image cone, while
a dense D-S3 interpolation map keeps the same section nonzero. The
obstruction concerns this coordinate-deletion criterion only.

\paragraph{I-33. Divisors on the one-dimensional torus and compatible weights.}
The prime-weight substitutions are integral cocharacters of the
split torus in D-S3. Pulling back the rational ratio of two products
gives divisors on the multiplicative line. Distinct prime root orbits
can overlap in at most one point, while the inseparable case is
controlled by integer divisor multiplicities. This forces any lost
ratio to be a character, which the coordinate variants separate.
The subset-sum example identifies the exact failure of closure under
products for the raw family. Interpolating its finite evaluation set
by one affine curve repairs scalar product closure because the
coordinate ring of the line is a domain. D-FGA places this curve and
the coefficient-image restriction; no purity theorem is needed.


\paragraph{G-18. Removing scalar orbit periods and comparing every large degree.}
The auxiliary Schur cone and its D-S6 Cartan coordinate modules remain
those of G-17. The smooth Fermat factor has finite projective stabilizer;
varying its padding line in a common stable open avoids every nontrivial
fixed locus. Unique factorization then gives a free projective orbit.
In the D-S3 central group scheme $\mu_4\subset\mathrm{SL}_4$, the
character exponent $ns-|\eta|=4\nu_4$ is trivial. Consequently two
distinct nonzero scalar multiples of the chosen tensor have distinct
closed orbits. Exact reductive invariants separate them, forcing the
degrees of common nonzero invariant sections to generate $\mathbb Z$.
The explicit numerical-semigroup argument gives sections in every
degree at least $\beta^2$. Multiplication by these sections is injective
in both quotient coordinate domains. D-FGA's standard Veronese and
Hilbert-function comparison can therefore be transported to every
integer degree, with the stated effective thresholds and common
asymptotic exponent. The rectangular central congruence is already
encoded by $ns=4\nu_4$. All stability, exact invariant-functor, and
quotient assertions here are in characteristic zero; an integral
model alone is not an assertion of the same behavior at every prime.

\paragraph{I-34. Local relation modules and coordinate cylinders.}
At one factor position, take the kernel of the univariate coefficient
map from the free module on product terms. A rank-one kernel with full
support is a circuit. The D-S3 prime cocharacters and the divisor
argument of I-33 force every canceled multiplicative ratio off this
position to be a character. A free coordinate separates its character
exponents, and the circuit property excludes every proper relation.
Higher local nullity is reduced by coefficient elimination. The result
gives the explicit bound on the number of free coordinate directions;
interpolation replaces the cylinders by integral cocharacters. In
D-FGA terms, the finite family preserves the coefficient-image module
after extension to the working field. The rank-four flattening example
defines a genuine nonzero tensor in the kernel of every one-coordinate
cylinder and identifies the precise failure of that smaller family.

\paragraph{P-14. Rank-one reset maps and scalar bridge sections.}
On a rank-one bundle-map chart, a nonzero pivot writes the map as a
column times a row divided by that pivot. Clearing all pivots splits
the scalar contraction into bridge sections. Invertible interior maps
and absorbed endpoint maps have identity constants over $F(t)$; the
ascending coefficient-row subspaces have at most $w-1$ strict growth
steps. One common interpolation morphism keeps every nonzero bridge
section nonzero, and their product remains nonzero in its coordinate
domain even when adjacent bridges share a reset variable. Constant
linear annihilators give exact coefficient-row preservation. This is
the D-FGA image-module criterion and the D-S3 dense substitution; no
equality of matrix image cones or purity statement is required. The
balanced rational-evaluation form follows from the stated partial
fractions and the same scalar-annihilator argument.


\paragraph{R-22. Cyclic characters and one open differential-rank condition.}
The permanent coefficient morphism is an integral polynomial map of
affine schemes. R-22 works over a cyclotomic extension in characteristic
zero. The D-S3 diagonal torus restricts to $\mu_m$, whose characters
split the differential coefficient space. A single staircase parameter
space supplies both local row and column cofactor relations. Their
highest components define two subspaces with explicitly computed
coprime-factor intersection. The D-S6 Schur determinant is the nonzero
minor of coefficient restriction; its principal geometric specialization
is a Vandermonde determinant with distinct cyclotomic exponents.
The product of finitely many nonzero parameter polynomials defines one
nonempty open subset on which all bidegrees are simultaneously valid.
An integer grid contains a point of that open after the root of unity
is fixed; translated motifs use the prescribed cyclotomic twists.
The final D-FGA differential minor is nonzero in one common deformation,
so its image dimension is at least the stated bound. The proof can be
spread over the localization where these particular denominators and
minors remain nonzero; it does not assert the same rank in every residue
characteristic. Its quotient and highest-weight consequences use the
characteristic-zero dictionaries of G-10 and G-17--G-18.

\paragraph{The G-17 quantitative corollary.}
The rank of a symmetric power and the six-factor Schur-rank polynomial
bound the ambient bundle rank in the D-S6 construction. Together with
the polynomial action degree, these give the binary-length estimate
for the existing first successful quotient degree. This is an effective
bound on the degrees of sections in the D-FGA graded algebra; it is
separate from computing the section or deciding its multiplicity.


\paragraph{R-23. A single graded part of a derived intersection.}
Let $Z=\operatorname{Spec}(R/J_m)$ be the permanent Jacobian scheme,
and let $L=\operatorname{Spec}S$ be a general four-plane in
$\operatorname{Spec}R$. Their D-Tor derived fiber product has
coordinate complex $(R/J_m)\otimes_R^{\mathbf L}S$.
Its first homology in internal degree $m$ is precisely the dimension
defect $\epsilon_m$. The known full-variable linear resolution gives
the preceding map, and the cyclic coefficient calculation proves
that its $2m-2$ old syzygies remain independent while no new constant
relation appears. The degree-$m$ differential then gives
$c_m=4m^2-2m+2-\epsilon_m$. This is an equality of dimensions of
explicit cohomology groups, not an inference from an alternating
$K_0$ class. The low-size positive defects prove failure of the
unqualified restriction argument. The depth bound rules out a full
regular-sequence justification but leaves the single graded
vanishing question open. D-FGA places the associated generic rank
stratum, and D-S6 distinguishes its actual homology from its
Euler characteristic.

\paragraph{P-15. A nonzero section killed by a dense interpolation morphism.}
The displayed quadratic is a rational coefficient tensor with
minimum tensor-train width three: its three balanced flattenings
have ranks $3,4,4$. The upper-triangular bundle-map factorization
keeps every interval on the invertibility open set. Nevertheless
its pullback under the specified centered Lagrange morphism is
zero. The identities in the three linear coordinates $S,U,V$
prove this directly; the harmonic cross-ratio identity supplies
the separately stated rational-evaluation examples. D-FGA records
the nonzero section and its zero pullback, while D-S3 records the
coordinate choices. The example is over $\mathbb Q(\mathrm i)$
for the fourth-root Lagrange nodes, and uses rational abscissas
for the other convention. No scalar change of coordinates is
silently omitted in comparing these two maps.

\subsection{What the manuscript record does and does not establish}
The supplied cote~133 study also reports decisions about the name on scan~008, a characteristic-two pre-$\lambda$ issue on scan~006, and a crossed-out Stiefel--Whitney assertion on scan~009. Those reports are retained as reports of earlier readers, with their own scope. No new visual decision about the name or the deleted text is made in this run. The five attached LaTeX transcriptions contain no original raster pages. The separate transcription table records the forced algebraic readings and the distinction between a textual ambiguity, a mathematical correction, and a claim whose ink cannot be checked here.
